\chapter[\emj{🌳}~Binary Search Trees (Árboles Binarios de Búsqueda)]{\emj{🌳}~Binary Search Trees (Árboles Binarios de Búsqueda)}

\begin{dsaquees}

Un árbol genealógico 🌳 de cabeza. La raíz (root) está hasta arriba.
Cada persona (nodo) puede tener máximo dos hijos.

Pero es un árbol inteligente con una regla de oro:
- Todo lo que está a la IZQUIERDA de un nodo es MENOR.
- Todo lo que está a la DERECHA de un nodo es MAYOR.

Es como una biblioteca 📚 donde los libros están tan bien ordenados
que puedes encontrar cualquiera en segundos, sin revisar todos los 
estantes.

\end{dsaquees}

\begin{dsaotraforma}

Un Binary Search Tree (BST) es una estructura de datos no lineal
formada por nodos. Cada nodo contiene un valor y punteros a sus hijos
izquierdo y derecho.

La propiedad de orden (izq < raíz < der) permite que las operaciones
de búsqueda, inserción y eliminación se realicen en tiempo 
logarítmico O(log n), siempre y cuando el árbol esté balanceado.

\end{dsaotraforma}

\begin{dsadiagrama}
\begin{verbatim}
BST de ejemplo:

         [8]          ← ROOT
       /     \
    [3]       [10]
   /   \         \
 [1]    [6]       [14]
       /   \      /
     [4]   [7]  [13]

- Leaf Nodes (hojas): 1, 4, 7, 13 (no tienen hijos).
- Height (altura): 3 (niveles desde la raíz).
\end{verbatim}
\end{dsadiagrama}

\begin{lstlisting}[language=C++]
1. Search: Comparas con la raíz. ¿Es menor? Vas a la izquierda. 
   ¿Es mayor? Vas a la derecha. Repites hasta encontrarlo o llegar 
   a un callejón sin salida (NULL).
2. Insert: Buscas el lugar donde debería estar el valor siguiendo la
   regla de orden y lo agregas como una nueva hoja.
3. Delete: Es la más compleja. Tiene 3 casos:
   - Nodo es una hoja: Solo lo borras.
   - Nodo tiene 1 hijo: El hijo sube a ocupar el lugar del padre.
   - Nodo tiene 2 hijos: Buscas al sucesor (el más pequeño del lado
     derecho) para que ocupe su lugar.
4. Traversal (Recorrido): 
   - Inorder: Izquierda -> Raíz -> Derecha (¡Te da los números en orden!).

📝 PSEUDOCÓDIGO
──────────────────────────────────────────────────────────────

// Búsqueda recursiva en un BST
función buscar(nodo, valor):
    SI nodo es NULL: return FALSO
    SI nodo.valor == valor: return VERDADERO
    SI valor < nodo.valor:
        return buscar(nodo.izq, valor)
    SINO:
        return buscar(nodo.der, valor)

💻 CÓDIGO C++
──────────────────────────────────────────────────────────────

#include <iostream>
using namespace std;

struct Node {
    int data;
    Node *left, *right;
    Node(int val) : data(val), left(nullptr), right(nullptr) {}
};

class BST {
public:
    Node* root;
    BST() : root(nullptr) {}

    // Insertar -> O(log n) promedio
    Node* insert(Node* node, int val) {
        if (!node) return new Node(val);
        if (val < node->data)
            node->left = insert(node->left, val);
        else
            node->right = insert(node->right, val);
        return node;
    }

    // Búsqueda -> O(log n) promedio
    bool search(Node* node, int val) {
        if (!node) return false;
        if (node->data == val) return true;
        if (val < node->data) return search(node->left, val);
        return search(node->right, val);
    }

    // Inorder Traversal -> O(n)
    void inorder(Node* node) {
        if (!node) return;
        inorder(node->left);
        cout << node->data << " ";
        inorder(node->right);
    }
};

int main() {
    BST tree;
    tree.root = tree.insert(tree.root, 8);
    tree.insert(tree.root, 3);
    tree.insert(tree.root, 10);
    tree.insert(tree.root, 1);
    tree.insert(tree.root, 6);

    cout << "Inorder: ";
    tree.inorder(tree.root); // 1 3 6 8 10 
    cout << endl;

    cout << "¿Está el 6? " << (tree.search(tree.root, 6) ? "Sí" : "No") << endl;

    return 0;
}

⏱️ COMPLEJIDAD (Big-O)
──────────────────────────────────────────────────────────────

  Operación   │ Promedio   │ Peor caso (Degenerado)
  ────────────┼────────────┼───────────────────────────
  Search      │ O(log n)   │ O(n)
  Insert      │ O(log n)   │ O(n)
  Delete      │ O(log n)   │ O(n)
  Traversal   │ O(n)       │ O(n)
  Espacio     │ O(n)       │ O(n)

* Peor caso ocurre cuando el árbol se vuelve una línea (como linked list).
\end{lstlisting}

\begin{dsaentrevistas}

✅ "Inorder Traversal = Sorted Array": Si recorres un BST en modo 
   Inorder, siempre obtendrás los elementos ordenados de menor a mayor.

💡 Lowest Common Ancestor (LCA): Un problema de Google súper común. 
   En un BST es fácil: el LCA es el primer nodo cuyo valor está ENTRE
   los dos valores que buscas.

⚠️ Balanceo: Un BST normal puede volverse muy ineficiente (O(n)) si 
   insertas los números ya ordenados. Para evitarlo existen los 
   Self-balancing Trees (como AVL).

🔑 Pregunta clave: "¿Es un Binary Tree o un Binary Search Tree?". 
   El segundo tiene la propiedad de orden, el primero no.

\end{dsaentrevistas}

\begin{dsaresumen}

• Izquierda menor, Derecha mayor.
• Búsqueda rápida O(log n) como adivinar un número.
• Inorder traversal = lista ordenada.
• El peor caso es O(n) si el árbol no está balanceado.
• Términos: Root (cima), Leaf (hoja), Parent (padre), Child (hijo).
• Altura del árbol = longitud del camino más largo a una hoja.
════════════════════════════════════════════════════════════════

\end{dsaresumen}


